% !TeX program = xelatex
\documentclass{article}
\usepackage[UTF8,fontset=fandol]{ctex}
\usepackage{amsmath,amssymb}
\defaultCJKfontfeatures{Script=CJK,Extension=.otf}
\setCJKmainfont[BoldFont=FandolSong-Bold,ItalicFont=FandolKai-Regular]{FandolSong-Regular}
\setCJKsansfont[BoldFont=FandolHei-Bold]{FandolHei-Regular}
\setCJKmonofont{FandolFang-Regular}
\newCJKfontfamily{\examplekai}[Scale=1.1]{FandolKai-Regular}
\newfontfamily{\examplelatin}[Numbers=OldStyle,Ligatures=TeX]{lmroman10-regular.otf}
\begin{document}
\section{中文与数学}
你好，世界！这是一个使用 ctex 和 Xe\LaTeX{} 的中文示例。
普通中文、English words 和行内公式 $f(x)=x^2+1$ 可以一起编辑。

\textbf{这是粗体中文。} \textit{这是楷体中文。}
\textsf{这是黑体中文。} \texttt{这是仿宋中文。}

{\heiti 黑体中文}、{\songti 宋体中文}、{\kaishu 楷体中文}、
{\fangsong 仿宋中文}都保留原来的西文字族。
{\zihao{-4}小四号字}与{\zihao{3}三号字}可以放在同一段中。
{\ziju{0.3}这几个字有额外的字间距。}

\noindent 这一段不缩进。前半句\hspace{1em}后半句之间有一个字宽的间隔。

\indent 这一段显式缩进。普通中文段落默认缩进两个汉字，
\verb|\noindent| 取消当前段落的首行缩进，\verb|\indent| 使用文档的缩进长度。
即时预览中的缩进标记可以点击修改，不会成为编译后的正文。

\begin{equation}\label{eq:chinese}
  \int_0^1 x^2\,dx=\frac{1}{3}.
\end{equation}
公式~\eqref{eq:chinese} 的标签也可以在即时预览中编辑。

\section{自动换行}
中文段落可以在汉字之间自动换行，并且标点符号不应出现在不合适的行首。
可以选择一段文字、输入新的汉字，或切换到源码窗口检查同一处的选择范围。

\section{字体与编辑练习}
在文档设置的 Fonts 页中选择字体或字体文件，并配置粗体、斜体和缩放。
这个示例使用 \TeX{} Live 自带的 Fandol 字体。
{\examplekai 这个命名字体使用楷体，并放大到原来的 1.1 倍。}
{\examplelatin 西文数字与连字示例：office AVATAR 0123456789，
\textsc{Small Capitals}，以及 ``quoted text''。}

\textsf{\textit{\textbf{嵌套的中文与 English 字体}}}，
\textbf{\textit{粗体与斜体组合}}，以及普通正文可以放在同一段中。
选择其中几个汉字，改变字族或字号，再清除格式。源码编辑器使用等宽字体，
中文字符使用可用的中文等宽字体或后备字体。

在公式前后输入中文：因此 $a>0$，也等价于 $E=mc^2$。
请尝试输入“不”“你好，世界！”，移动光标到段尾，并在即时视图与完整源码之间
比较选择范围。按 Ctrl+S 保存，再打开编译后的 PDF 比较字体与换行。
\end{document}
